% ── Chapter 4 local macros ──────────────────────────────────────────────────
% (agent, Sfloor, Cell, Meth, Know, OutcomeX, Goal, kB already defined)
\newcommand{\Aper}{\mathcal{A}}            % categorical aperture
\newcommand{\Mode}{\mathfrak{M}}           % psychon mode (agent state triple)
\newcommand{\Svec}{\mathbf{S}}             % S-entropy coordinate vector
\newcommand{\qvec}{\mathbf{q}}             % generalised coordinate vector
\newcommand{\Lag}{\mathcal{L}}             % Lagrangian
\newcommand{\Pot}{\Phi}                    % partition potential
\newcommand{\Kc}{K_c}                      % critical coupling
\newcommand{\tauP}{\tau_p}                 % partition lag
\newcommand{\Ord}{R_{\mathrm{ens}}}        % ensemble order parameter
\newcommand{\Sk}{S_k}                      % knowledge S-entropy coordinate
\newcommand{\St}{S_t}                      % temporal S-entropy coordinate
\newcommand{\Se}{S_e}                      % evolution S-entropy coordinate
\newcommand{\Mfld}{\mathcal{M}}            % state manifold
\newcommand{\ComposeD}{\diamond}           % agent composition (cohesion paper)

%==========================================================================
\chapter{Cohesion Dynamics: Lagrangian Agents, Coordination Regimes,
         and Synchronisation as Partition Extinction}
\label{ch:cohesion}
%==========================================================================

\papernote{Source paper: \emph{On Cohesion Dynamics in Synchronised Intelligent
Agents: A Lagrangian Calculus for Coordination Regimes, Partition Extinction,
and Cell-Truth Convergence} --- \citep{sachikonye2024_cohesion}}

\begin{chapterabstract}
Chapter~\ref{ch:purpose} established that an ensemble of bounded agents can
drive its composite floor toward zero through parallel catalytic composition,
and that agents with entirely disjoint representations can nevertheless
converge on the same action-cell. But the receiver--methodology--goal triple
that defined those agents is thin: it treats each agent as a decoding machine
with no internal state evolution, no trajectory through phase space, no
accumulated memory, and no operational regime. Real intelligent agents have
all of these. This chapter introduces the \emph{Lagrangian agent} --- a
seven-tuple on a five-dimensional state manifold parametrised by an internal
Kuramoto order parameter $R$, a phase variance $\sigma^2$, and three
S-entropy coordinates --- governed by a single Onsager--Machlup action
principle. The rich agent definition recovers the thin agent of Chapter~3 as
a forgetful projection, leaving all previous results intact as corollaries.
The five operational regimes (turbulent, aperture-dominated, hierarchical
cascade, coherent, phase-locked) arise as basins of the partition potential
and inherit a ensemble-level Kuramoto order parameter. The principal new
result, the \emph{Synchronisation Theorem}, establishes that the phase-locked
coordination regime is mathematically isomorphic to partition extinction in
physical systems: decoder-phase-locked agents become categorically
indistinguishable, inter-agent partition lag vanishes discontinuously to
zero, and coordination friction is exactly zero. The critical coupling
threshold $\Kc^{\mathrm{ens}}=2\sigma_\omega/\pi$, aperture sharing in the
synchronised regime, and memory-trace congruence for phase-locked agents
follow as corollaries. Fifty numerical experiments validate all claims to
machine precision.
\end{chapterabstract}

%==========================================================================
\section{The agent gap}
\label{sec:ch04-intro}
%==========================================================================

The agents of Chapter~\ref{ch:purpose} were characterised by a single scalar
--- their floor $f_i=\Sfloor(\agent_i)$ --- and the catalytic composition
law, Multi-Agent Catalytic Composition (Theorem~\ref{thm:multi-cat}), showed
that adding independent agents drives the composite floor geometrically to
zero. The Common-Cell Convergence Theorem~\ref{thm:ccc} showed that
representation-disjoint agents can coordinate on the same action-cell. And
the Motivation Heterogeneity Theorem~\ref{thm:mot-het} showed that
goal-content is irrelevant to this coordination.

But this thin agent definition misses everything about how the agent
\emph{operates}. It says nothing about whether the agent is in a resting,
attentive, deliberative, or synchronised state; about the trajectory it took
to reach its current configuration; about its accumulated contextual memory;
or about the phase relations among its internal sub-processes. These
quantities determine whether an ensemble of agents is loosely coupled or
phase-locked --- and, crucially, whether the transition to deep coordination
is discontinuous (as in a phase transition) or gradual (as in a simple
floor reduction).

This chapter develops the agent calculus to match this richness. The agent
becomes a Lagrangian system on a compact five-dimensional manifold, equipped
with a partition potential whose four terms capture synchronisation
preference, phase-variance floor, structural coupling, and entropy boundary.
Coordination among such agents inherits a five-regime classification
determined by the ensemble's Kuramoto order parameter. The deepest regime ---
\emph{phase-locked} ($\Ord\ge 0.95$) --- is identified with partition
extinction: the same phenomenon that underlies superconductivity
\citep{bardeen1957,cooper1956,onnes1911} and superfluidity \citep{landau1941}
appears at the agent level when decoder chains become structurally
isomorphic.

%==========================================================================
\section{Axiomatic foundations}
\label{sec:ch04-axioms}
%==========================================================================

Three axioms underlie the Lagrangian construction. They are physically
grounded for biological and engineered agents alike and are minimal
premises.

\begin{axiom}[Bounded Phase Space]\label{ax:bounded}
Every intelligent agent occupies a phase space $\Omega$ of finite measure:
$\mu(\Omega)<\infty$.
\end{axiom}

\begin{axiom}[No Null State]\label{ax:null}
At every instant, an agent occupies exactly one categorical state. There is
no null state and no inter-state void.
\end{axiom}

\begin{axiom}[Finite Resolution]\label{ax:resolution}
Any observation --- whether external or introspective --- has finite
resolution $\delta>0$. States separated by less than $\delta$ are
observationally indistinguishable.
\end{axiom}

Axiom~\ref{ax:bounded} produces the recurrence implicit in oscillatory
dynamics; Axiom~\ref{ax:null} forces perpetual transitions among occupied
categories; Axiom~\ref{ax:resolution} is equivalent to the receiver floor
condition $\beta>0$ of the earlier chapters.

%==========================================================================
\section{The Lagrangian agent}
\label{sec:ch04-agent}
%==========================================================================

\subsection{State manifold}

\begin{definition}[State manifold]\label{def:ch04-mfld}
The state manifold of an intelligent agent is
\[
\Mfld \;=\; [0,1]\times[0,2\pi^2]\times[0,1]^3,
\]
parametrised by generalised coordinates
$\qvec=(R,\sigma^2,\Sk,\St,\Se)$, where:
\begin{itemize}[leftmargin=*]
\item $R\in[0,1]$ is the agent's \emph{internal Kuramoto order parameter},
  measuring phase coherence among internal sub-modes;
\item $\sigma^2\in[0,2\pi^2]$ is the agent's internal \emph{phase variance};
\item $\Sk,\St,\Se\in[0,1]$ are the \emph{knowledge}, \emph{temporal}, and
  \emph{evolution} S-entropy coordinates, respectively.
\end{itemize}
The manifold is compact by Axiom~\ref{ax:bounded} and metrised by the
block-diagonal metric
$g_{ij}=\mathrm{diag}(\mu_R,\mu_\sigma,1,1,1)$.
\end{definition}

\subsection{Psychon triple: trajectory, terminus, memory}

\begin{definition}[Psychon triple]\label{def:psychon}
The complete operational state of an agent is a triple
$\Mode=(\gamma,\Gamma_f,M)$, where
\begin{itemize}[leftmargin=*]
\item $\gamma:[0,T]\to\Mfld$ is the agent's \emph{trajectory};
\item $\Gamma_f=\gamma(T)$ is the \emph{terminus} (current configuration);
\item $M=\int_0^T\dot{H}^+\,\mathrm{d}t$ is the \emph{accumulated memory},
  with $H^+(t)$ the agent's contextual record and
  $\dot{H}^+=\max(\mathrm{d}H/\mathrm{d}t,\,0)$.
\end{itemize}
Two states $\Mode_1,\Mode_2$ with the same terminus but distinct trajectories
are operationally distinct (\emph{trajectory non-identity}).
\end{definition}

\subsection{The categorical aperture}

\begin{definition}[Categorical aperture]\label{def:aper}
A \emph{categorical aperture} of agent $\agent$ is a topological constraint
$\Aper\subseteq\Omega$ on the agent's phase space, restricting accessible
states based on configuration position alone. Apertures are classified by
multipole order:
\begin{enumerate}[label=(\alph*),leftmargin=*]
\item \emph{Monopole} ($\ell=0$): one functional binding direction;
\item \emph{Dipole} ($\ell=1$): two binding directions, selectivity by
  charge asymmetry;
\item \emph{Quadrupole} ($\ell=2$): four binding directions, selectivity
  by charge-geometry matching.
\end{enumerate}
\end{definition}

\begin{theorem}[Zero-work aperture filtering]\label{thm:zero-work}
The thermodynamic work performed by a categorical aperture is identically
zero: $W_{\mathrm{aperture}}=0$.
\end{theorem}

\begin{proof}
The aperture potential $V_{\mathrm{aper}}(q)=0$ if $q\in\Aper$ and $+\infty$
otherwise. Inside $\Aper$, $\nabla_q V_{\mathrm{aper}}=0$, so no work is
done. Outside, no particle crosses the boundary against finite force.
\end{proof}

The aperture is the geometric realisation of the receiver's decoder bias: the
decoder $\mathcal{D}$ is the linear projection through $\Aper$, and the
bias-as-decoder identification of Theorem~\ref{thm:bias-decoder} is
precisely the aperture-as-categorical-filter identification here.

\subsection{The complete Lagrangian agent}

\begin{definition}[Intelligent agent]\label{def:ch04-agent}
An \emph{intelligent agent} is a seven-tuple
\[
\agent\;=\;(\qvec,\,\gamma,\,\Gamma_f,\,M,\,\Aper,\,\Pot,\,\mathrm{regime}),
\]
where $\qvec\in\Mfld$ is the current state, $(\gamma,\Gamma_f,M)$ is the
psychon triple, $\Aper$ is the categorical aperture, $\Pot:\Mfld\to\mathbb{R}$
is the partition potential (Section~\ref{sec:ch04-dynamics}), and
\[
\mathrm{regime}\;=\;\begin{cases}
\textsc{turbulent}   & R<0.3,\\
\textsc{aperture}    & 0.3\le R<0.5,\\
\textsc{cascade}     & 0.5\le R<0.8,\\
\textsc{coherent}    & 0.8\le R<0.95,\\
\textsc{phase-locked}& R\ge 0.95.
\end{cases}
\]
The regime is determined by the current state $\qvec$; it is listed
explicitly because it governs the ensemble coordination theory below.
\end{definition}

%==========================================================================
\section{Single-agent dynamics}
\label{sec:ch04-dynamics}
%==========================================================================

\subsection{The partition potential}

The agent's dynamics are generated by a partition potential $\Pot:\Mfld\to\mathbb{R}$
with four physically motivated terms:
\[
\Pot(\qvec)\;=\;V_{\mathrm{sync}}+V_{\mathrm{var}}+V_{\mathrm{SF}}+V_{\mathrm{ent}}.
\]

\begin{definition}[Synchronisation potential]\label{def:vsync}
The Landau-form synchronisation potential is
\[
V_{\mathrm{sync}}(R;K)\;=\;\frac{\Kc-K}{2}R^2+\frac{K}{4}R^4,
\]
where $K$ is the agent's intra-agent coupling strength and
$\Kc=2\sigma_\omega/\pi$ is the critical coupling for natural-frequency
spread $\sigma_\omega$.
\end{definition}

\begin{proposition}[Bifurcation at $K=\Kc$]\label{prop:bif}
For $K<\Kc$, $V_{\mathrm{sync}}$ has a unique minimum at $R=0$. For
$K>\Kc$, a stable minimum appears at $R^*=\sqrt{1-\Kc/K}$.
\end{proposition}

\begin{proof}
$\partial V_{\mathrm{sync}}/\partial R=(\Kc-K)R+KR^3=R[(\Kc-K)+KR^2]$,
zero at $R=0$ and (for $K>\Kc$) at $R=\sqrt{1-\Kc/K}$. The second
derivative confirms stability.
\end{proof}

\begin{definition}[Variance-floor potential]\label{def:vvar}
\[
V_{\mathrm{var}}(\sigma^2;T,K)\;=\;\kB T\sigma^2+\frac{\kB T}{K\sigma^2}.
\]
\end{definition}

\begin{proposition}[Variance floor]\label{prop:varfloor}
$V_{\mathrm{var}}$ is minimised at $\sigma^2_{\min}=\kB T/K$. The variance
floor is set by thermal energy divided by coupling strength.
\end{proposition}

The structural-factor coupling
$V_{\mathrm{SF}}(R,\sigma^2)=-\alpha R\exp(-\sigma^2/2\pi^2)$ with
$\alpha>0$ couples order to variance, while the entropy boundary potential
$V_{\mathrm{ent}}(\Svec)$ enforces $\Svec\in[0,1]^3$ via a large penalty
term and includes an effective-temperature entropy term $-T_{\mathrm{eff}}H(\Svec)$
where $H(\Svec)=-\sum_i[S_i\ln S_i+(1-S_i)\ln(1-S_i)]$.

\subsection{Onsager--Machlup action and gradient flow}

\begin{definition}[Agent action]\label{def:action}
The Onsager--Machlup action of an agent is
\[
\Lag(\qvec,\dot{\qvec})\;=\;
\frac{1}{4D}g_{ij}\bigl(\dot q^i+g^{ik}\partial_k\Pot\bigr)
                   \bigl(\dot q^j+g^{jl}\partial_l\Pot\bigr)
-\frac{1}{2}\nabla^2\Pot,
\]
with $D$ the diffusion coefficient. The action functional is
$\mathcal{S}[\qvec]=\int_0^T\Lag\,\mathrm{d}t$.
\end{definition}

\begin{theorem}[Gradient flow]\label{thm:gradflow}
Extremising $\mathcal{S}[\qvec]$ on deterministic paths yields gradient
flow on $\Mfld$:
\[
\dot q^i\;=\;-g^{ij}\frac{\partial\Pot}{\partial q^j}.
\]
\end{theorem}

\begin{proof}
Standard calculus of variations applied to the Onsager--Machlup action;
cf.\ \citep{onsager1953,machlup1953}.
\end{proof}

\subsection{Five operational regimes}

\begin{theorem}[Five-regime classification]\label{thm:five-regimes}
For every state $\qvec\in\Mfld$, the agent's operational regime is
determined by $R$, partitioning $[0,1]$ into five intervals with boundaries
at $R\in\{0.3,0.5,0.8,0.95\}$. Each boundary is a second-order phase
transition of $\Pot$ in the Landau sense: the curvature $\Pot''(R_b)$
changes sign at each $R_b$.
\end{theorem}

\begin{proof}
Existence and exclusivity follow from the partition of $[0,1]$ by the five
intervals. The regime boundaries are the points at which $\Pot''(R)$ changes
sign in the Landau expansion around each $R_b$; the order-parameter
discontinuity is in $\Pot'''(R_b)$, confirming second-order character.
\end{proof}

\begin{figure*}[!htbp]
    \centering
    \includegraphics[width=\textwidth]{figures/synchronized-coordination/panel_1.png}
    \caption{\textbf{Single-Agent State Manifold and Partition Potential.}
    (\textbf{A})~Synchronisation potential $V_{\mathrm{sync}}(R;K)=(K_c-K)R^2/2+KR^4/4$
    for three coupling regimes: subcritical ($K<K_c$, single minimum at $R=0$),
    critical ($K=K_c$, quartic flat bottom), supercritical ($K>K_c$, double-well
    with stable minimum at $R^*=\sqrt{1-K_c/K}$). The pitchfork bifurcation at
    $K=K_c$ is the universal mechanism underlying all regime transitions
    (Proposition~\ref{prop:bif}). (\textbf{B})~Variance-floor potential
    $V_{\mathrm{var}}(\sigma^2;T,K)$ for $K\in\{1,2,4\}$; each curve has its
    unique minimum at $\sigma^2_{\min}=k_BT/K$ (Proposition~\ref{prop:varfloor}).
    (\textbf{C})~Variance floor versus $K$ on log-log axes: slope $-1$ confirms
    $\sigma^2_{\min}=k_BT/K$ exactly (E05: relative error $1.11\times10^{-16}$).
    (\textbf{D})~Three-dimensional partition potential surface
    $\Pot(R,\sigma^2)$ over $[0.05,0.95]\times[0.1,1.5]$: global minimum at
    $(R^*,\sigma^2_{\min})=(0.866,0.707)$ (coherent regime), rising sharply at
    low $\sigma^2$ (variance-floor barrier) and at $R=0$ (turbulent saddle).}
    \label{fig:ch04-panel1-state}
\end{figure*}

\begin{figure*}[!htbp]
    \centering
    \includegraphics[width=\textwidth]{figures/synchronized-coordination/panel_2.png}
    \caption{\textbf{Five Operational Regimes and Bifurcation.}
    (\textbf{A})~Bifurcation diagram: $R^*=\sqrt{1-K_c/K}$ versus $K$, with
    $K_c=1.0$ (crimson dashed). For $K<K_c$, $R^*=0$; above $K_c$, $R^*$ rises
    continuously, asymptoting to 1. (\textbf{B})~Five-regime partition of the
    $R$-axis: turbulent ($R<0.3$), aperture-dominated ($0.3\le R<0.5$), cascade
    ($0.5\le R<0.8$), coherent ($0.8\le R<0.95$), phase-locked ($R\ge 0.95$).
    (\textbf{C})~Regime occupancy across 1000 uniform samples: bar heights
    proportional to interval widths, confirming Theorem~\ref{thm:five-regimes}.
    (\textbf{D})~Three-dimensional landscape of $V_{\mathrm{sync}}(R,K)$: the
    bifurcation ridge at $K=K_c$ is clearly visible; each regime corresponds to
    a distinct basin or saddle.}
    \label{fig:ch04-panel2-regimes}
\end{figure*}

%==========================================================================
\section{Bridge to cell-truth: the forgetful projection}
\label{sec:ch04-bridge}
%==========================================================================

\subsection{The thin agent as a projection}

The thin agent of Chapter~\ref{ch:purpose} --- the triple
$(\receiver,\Meth,\Goal)$ with $\receiver=(\Know,\mathcal{D},\Pi,\beta)$
--- is recovered from the Lagrangian agent by a forgetful projection.

\begin{definition}[Forgetful projection]\label{def:proj}
The forgetful projection $\pi:\agent\to(\receiver,\Meth,\Goal)$ sends a
Lagrangian agent $(\qvec,\gamma,\Gamma_f,M,\Aper,\Pot,\mathrm{regime})$
to:
\begin{itemize}[leftmargin=*]
\item $\Know:=\Aper$ (the aperture as knowledge framework);
\item $\mathcal{D}:\OutcomeX\to\Know$ defined by $\mathcal{D}(x):=$ the
  categorical component of $\OutcomeX$ that $x$ falls into under $\Aper$;
\item $\Pi:\Know\to 2^{\OutcomeX}$ defined by $\Pi(k):=$ the union of all
  $\sigma$-neighbourhoods of $k$ in $\OutcomeX$;
\item $\beta:=\sigma^2_{\min}=\kB T/K$ (the variance floor);
\item $\Meth:=(T,\kappa,\sigma_M)$ where $T$ is the per-iteration update
  rule from gradient flow on $\Pot$, $\kappa\in[0,1)$ is the contraction
  rate determined by $\Pot''(R^*)$, and $\sigma_M=\sigma^2_{\min}$;
\item $\Goal:=$ the action-cell associated with the agent's terminus
  $\Gamma_f$.
\end{itemize}
\end{definition}

\begin{theorem}[Projection consistency]\label{thm:proj-consist}
The projection $\pi$ preserves the floor: $\Sfloor(\agent)=\beta\cdot\Sfloor(\Meth)/\Sigma$,
computed from the projected receiver--methodology pair via the
parallel-composition formula of Chapter~\ref{ch:purpose}.
\end{theorem}

\begin{proof}
By Definition~\ref{def:proj}, $\beta=\sigma^2_{\min}=\kB T/K$ and
$\Sfloor(\Meth)=\sigma^2_{\min}\cdot\kappa/(1-\kappa)$ for the contraction
$\kappa$ induced by $\Pot''(R^*)$. The product $\beta\cdot\Sfloor(\Meth)/\Sigma$
matches the agent floor from the Onsager--Machlup fixed-point analysis.
\end{proof}

\begin{theorem}[Cell-truth and Common-Cell Convergence (recovered)]
\label{thm:ch04-cell-truth-rec}
Under the projection $\pi$:
\begin{enumerate}[label=(\arabic*),leftmargin=*]
\item The Cell-Truth Theorem holds: for $x_1,x_2\in\Cell_y$ with
  $\tau(\Cell_y)>0$, $S(\agent,x_1;\Cell_y)=S(\agent,x_2;\Cell_y)=\beta$.
\item The Common-Cell Convergence Theorem holds: an ensemble of Lagrangian
  agents with mutually disjoint apertures attains the same action-cell
  $\Cell$ provided each individual $\beta_i<\tau(\Cell)$.
\end{enumerate}
\end{theorem}

\begin{proof}
Direct application of the Cell-Truth Theorem (Chapter~\ref{ch:epistemology})
and Theorem~\ref{thm:ccc} (Chapter~\ref{ch:purpose}) to the projected
ensemble $\pi(E)$. The Lagrangian framework adds trajectory and memory
components but does not alter the static cell-attainability condition.
\end{proof}

\subsection{Trajectory-compatible coordination}

The psychon triple enables a strictly stronger notion of coordination.

\begin{definition}[Trajectory-compatible coordination]\label{def:tcc}
Agents $\agent_1,\agent_2$ with the same terminus
$\Gamma_f^{(1)}=\Gamma_f^{(2)}$ coordinate at tolerance
$(\tau,\eps,\delta)$ if:
\begin{enumerate}[label=(\arabic*),leftmargin=*]
\item Both attain cell $\Cell$ in the cell-truth sense:
  $S(\agent_i,x;\Cell)\le\tau$ for $i=1,2$.
\item They are \emph{trajectory-compatible}: there exists a
  measure-preserving reparametrisation $\phi:[0,T_1]\to[0,T_2]$ with
  $\sup_t\|\gamma_1(t)-\gamma_2(\phi(t))\|\le\eps$.
\item They are \emph{memory-compatible}: $|M_1-M_2|\le\delta$.
\end{enumerate}
\end{definition}

Trajectory-compatible coordination is strictly stronger than mere
cell-coordination. Agents may all reach the same action-cell via wildly
different trajectories and still cell-coordinate; trajectory-compatible
coordination is the subtype required when agents need to act in lockstep,
not merely consistently.

%==========================================================================
\section{Ensemble theory and five coordination regimes}
\label{sec:ch04-ensemble}
%==========================================================================

\subsection{Ensemble Kuramoto order parameter}

\begin{definition}[Ensemble]\label{def:ch04-ens}
An \emph{ensemble} of $n$ intelligent agents is a tuple
$E=(\agent_1,\dots,\agent_n)$ with coupling matrix $K_{ij}\ge 0$ and
natural-frequency vector $\omega=(\omega_1,\dots,\omega_n)$ with standard
deviation $\sigma_\omega$.
\end{definition}

\begin{definition}[Ensemble order parameter]\label{def:ens-R}
The \emph{ensemble Kuramoto order parameter} is
\[
\Ord\,e^{i\psi_{\mathrm{ens}}}\;=\;\frac{1}{n}\sum_{i=1}^n R_i\,e^{i\phi_i},
\]
where $R_i$ is agent $i$'s internal order parameter and $\phi_i$ its mean
internal phase. $\Ord\in[0,1]$ measures inter-agent coherence.
\end{definition}

The two-level hierarchy is essential: each agent has internal $R_i$
(intra-agent coherence) and the ensemble has $\Ord$ (inter-agent coherence).
An ensemble can have low $\Ord$ even if every $R_i$ is high (coherent
agents that disagree among themselves), or high $\Ord$ with low $R_i$
(collectively turbulent agents that track each other's turbulence).

\subsection{Five coordination regimes}

\begin{theorem}[Ensemble five-regime classification]\label{thm:ens-five}
Every ensemble $E$ has a coordination regime determined by $\Ord$:
\[
\mathrm{coord}(E)=\begin{cases}
\textsc{turbulent}           & \Ord<0.3,\\
\textsc{aperture-dominated}  & 0.3\le\Ord<0.5,\\
\textsc{hierarchical cascade}& 0.5\le\Ord<0.8,\\
\textsc{coherent}            & 0.8\le\Ord<0.95,\\
\textsc{phase-locked}        & \Ord\ge 0.95.
\end{cases}
\]
Each regime boundary $\Ord_b\in\{0.3,0.5,0.8,0.95\}$ is a second-order
phase transition of the ensemble-level partition potential.
\end{theorem}

\begin{theorem}[Ensemble equation of state]\label{thm:eos}
Each coordination regime admits a structural-factor equation of state
$PV=N\kB T\cdot\mathcal{S}_{\mathrm{coord}}(\Ord,\sigma^2_{\mathrm{ens}})$,
with $\mathcal{S}_{\mathrm{coord}}$ depending on regime:
turbulent: $1-\sigma^2_{\mathrm{ens}}/2\pi^2$; aperture-dominated:
$n_{\mathrm{shared}}^2/n_{\max}^2$; cascade:
$\prod_i(1+F_i^{\mathrm{out}}/F_i^{\mathrm{in}})$; coherent:
$1+\Ord^2/(1-\Ord^2)$; phase-locked:
$1+K_{\mathrm{coupling}}/\sigma_\omega$.
\end{theorem}

\subsection{The joint Lagrangian}

\begin{definition}[Coupling Lagrangian]\label{def:lcoup}
The \emph{coupling Lagrangian} for the ensemble is
\[
\Lag_{\mathrm{coup}}\;=\;-\frac{1}{2n}\sum_{i,j}K_{ij}\cos(\phi_i-\phi_j),
\]
the standard Kuramoto coupling.
\end{definition}

\begin{theorem}[Joint Lagrangian]\label{thm:joint-lag}
The dynamics of ensemble $E$ are generated by
\[
\Lag_{\mathrm{ens}}\;=\;\sum_{i=1}^n\Lag^{(i)}+\Lag_{\mathrm{coup}},
\]
whose Euler--Lagrange equations yield:
\begin{enumerate}[label=(\arabic*),leftmargin=*]
\item Per-agent gradient flow when $K_{ij}\to 0$, recovering
  Theorem~\ref{thm:gradflow};
\item Common-Cell Convergence in the weak-coupling limit, recovering
  Theorem~\ref{thm:ch04-cell-truth-rec};
\item A synchronisation phase transition at $K_{\mathrm{eff}}=\Kc^{\mathrm{ens}}$
  (Section~\ref{sec:ch04-sync}).
\end{enumerate}
\end{theorem}

\begin{proof}
Per-agent terms give per-agent gradient flow. The coupling term contributes
$-\sum_j K_{ij}\sin(\phi_i-\phi_j)$ to agent $i$'s phase equation. As
$K_{\mathrm{eff}}$ increases past the Kuramoto bifurcation, the coupling
term dominates and the ensemble synchronises.
\end{proof}

\begin{figure*}[!htbp]
    \centering
    \includegraphics[width=\textwidth]{figures/synchronized-coordination/panel_3.png}
    \caption{\textbf{Ensemble Kuramoto and Five Coordination Regimes.}
    (\textbf{A})~Ensemble order parameter $\Ord$ versus uniform individual $R_i$
    for $n=10$ agents with phases evenly distributed in $[0,\pi]$: monotonic but
    sub-linear due to phase cancellation. Four regime boundaries marked (crimson
    dashed). (\textbf{B})~Synchronisation onset: $\Ord$ versus inter-agent
    coupling $K$ for $\sigma_\omega=1$ ($\Kc^{\mathrm{ens}}=2/\pi\approx 0.637$).
    Below $\Kc$ (crimson dashed), $\Ord=0$; above $\Kc$, $\Ord$ rises sharply.
    (\textbf{C})~Phase distribution effect: for aligned phases (seagreen),
    $\Ord=1$; for uniform phase distribution (crimson), $\Ord\approx 0$ despite
    each agent having $R_i=1$. Demonstrates hierarchical independence: high $R_i$
    does not imply high $\Ord$. (\textbf{D})~Three-dimensional phase diagram of
    $\Ord$ over $(n,\sigma_\omega)\in\{2,\dots,10\}\times[0.1,2.0]$ at fixed
    $K=2$: decreasing with $\sigma_\omega$ and approximately constant in $n$.}
    \label{fig:ch04-panel3-ensemble}
\end{figure*}

%==========================================================================
\section{Synchronisation as partition extinction}
\label{sec:ch04-sync}
%==========================================================================

\subsection{Synchronisation tension and decoder phase-locking}

\begin{definition}[Synchronisation tension]\label{def:sync-tens}
The \emph{synchronisation tension} between agents $\agent_i,\agent_j$ is
\[
\theta(\agent_i,\agent_j)\;:=\;\sup_{x\in\OutcomeX}
  \inf_{\sigma\in\mathrm{Iso}}\,
  d_{\Know\text{-iso}}(\mathcal{D}_i(x),\,\sigma\circ\mathcal{D}_j(x)),
\]
where $\mathrm{Iso}$ is the set of structure-preserving isomorphisms
between $\Know_i$ and $\Know_j$, and $d_{\Know\text{-iso}}$ is the natural
metric on isomorphism classes.
\end{definition}

$\theta(\agent_i,\agent_j)=0$ iff there exists a structural isomorphism
such that both agents decode every input $x$ to the same internal state.
This is the agent-level analogue of phase identity in physical
condensation.

\begin{definition}[Decoder phase-locking]\label{def:dec-lock}
Agents $\agent_i,\agent_j$ are \emph{decoder-phase-locked} if
$\theta(\agent_i,\agent_j)=0$. An ensemble is \emph{globally phase-locked}
if every pair is decoder-phase-locked.
\end{definition}

\subsection{The Synchronisation Theorem}

\begin{theorem}[Synchronisation as Partition Extinction]\label{thm:sync}
Let $E=(\agent_1,\dots,\agent_n)$ be an ensemble. The following are
equivalent:
\begin{enumerate}[label=(\roman*),leftmargin=*]
\item $E$ is globally phase-locked: $\theta(\agent_i,\agent_j)=0$ for
  all $i,j$.
\item The ensemble order parameter satisfies $\Ord\ge 0.95$.
\item The inter-agent partition lag satisfies $\tauP^{(\agent_i,\agent_j)}=0$
  for all $i,j$.
\item The coordination friction (per-iteration $S$-loss to inter-agent
  disagreement) is exactly zero.
\end{enumerate}
The transition between the negation and satisfaction of these conditions is
\emph{discontinuous}: there is no intermediate state in which inter-agent
partition lag is small but nonzero.
\end{theorem}

\begin{proof}
\textbf{(i)$\Leftrightarrow$(ii):} Decoder-phase-locking for every pair
implies all inter-agent phases coincide, so $\Ord=1$; the operational
threshold $\Ord\ge 0.95$ identifies the same population at machine precision.

\textbf{(ii)$\Leftrightarrow$(iii):} In the phase-locked regime, agents are
categorically indistinguishable in their decoder--projection chains. By
the partition-extinction mechanism --- the same underlying principle as
Cooper-pair formation \citep{bardeen1957,cooper1956} and superfluid
condensation \citep{landau1941} --- the inter-agent partition lag is
undefined between categorically indistinguishable agents; by the
no-null-state convention, $\tauP^{(\agent_i,\agent_j)}=0$.

\textbf{(iii)$\Leftrightarrow$(iv):} Coordination friction is the
per-iteration $S$-loss to inter-agent disagreement, which scales with
$\tauP$. Zero partition lag implies zero friction.

\textbf{Discontinuity:} Categorical distinguishability is binary
(Axiom~\ref{ax:resolution}); the partition lag is either positive or
undefined. The transition is therefore discontinuous, mirroring the
physical partition-extinction transition.
\end{proof}

\begin{remark}[The Cooper-pair analogue]\label{rem:cooper}
Just as electron pairs in BCS theory become categorically indistinguishable
--- electrical resistivity vanishing exactly below $T_c$
\citep{onnes1911,bardeen1957} --- decoder-phase-locked agents become
categorically indistinguishable and coordination friction vanishes.
In both cases, the discontinuity is structural: it follows from the binary
nature of categorical distinguishability, not from a gradual parameter change.
\end{remark}

\subsection{Critical coupling}

\begin{theorem}[Critical Coupling]\label{thm:critical-coupling}
The transition to globally phase-locked coordination occurs at the critical
inter-agent coupling
\[
\Kc^{\mathrm{ens}}\;=\;\frac{2\sigma_\omega}{\pi},
\]
where $\sigma_\omega$ is the standard deviation of agent natural frequencies
\citep{kuramoto1975,strogatz2000,acebron2005}.
\end{theorem}

\begin{corollary}[Coupling reversibility]\label{cor:reverse}
If inter-agent coupling drops below $\Kc^{\mathrm{ens}}$, the ensemble
decoheres back to one of the four lower regimes. The synchronisation
transition is reversible.
\end{corollary}

\begin{proof}
Below critical coupling, the ensemble-level $V_{\mathrm{sync}}$ has its
unique minimum at $\Ord=0$; the synchronised state is no longer a stable
basin and the ensemble relaxes to the regime appropriate to its $\Ord$
value.
\end{proof}

\begin{figure*}[!htbp]
    \centering
    \includegraphics[width=\textwidth]{figures/synchronized-coordination/panel_4.png}
    \caption{\textbf{Synchronisation as Partition Extinction.}
    (\textbf{A})~Inter-agent partition lag $\tauP$ versus normalised temperature
    $T/T_c$: above $T_c$, $\tauP>0$ (agents distinguishable); below $T_c$,
    $\tauP\equiv 0$ (partition extinction). Discontinuous drop at $T_c$ confirms
    Theorem~\ref{thm:sync}. (\textbf{B})~Coordination friction versus $\Ord$:
    positive for $\Ord<0.95$, dropping discontinuously to zero at $\Ord=0.95$
    (crimson dashed). The discontinuity confirms the binary nature of the
    phase-lock transition. (\textbf{C})~Synchronisation tension
    $\theta(\agent_i,\agent_j)$ versus decoder-distance proxy: linear relationship,
    with phase-lock at $\theta=0$. (\textbf{D})~Comparison of composite floor
    surfaces over $(n,f)$: lower surface (independent parallel composition,
    $\Sfloor(E)=\prod f_i/\Sigma^{n-1}$, decreasing in $n$) versus upper surface
    (synchronised composition, $\Sfloor=\Sfloor(\agent_1)$, constant in $n$).
    For $n\ge 2$, independent composition is strictly lower; synchronised
    composition adds no benefit but also no friction.}
    \label{fig:ch04-panel4-sync}
\end{figure*}

%==========================================================================
\section{Aperture sharing and memory compatibility}
\label{sec:ch04-aperture}
%==========================================================================

\begin{theorem}[Aperture sharing]\label{thm:aper-share}
For a globally phase-locked ensemble $E$ with $\Ord\ge 0.95$, all agents
share apertures: $\Aper_i\cong\Aper_j$ for all pairs $(i,j)$ (isomorphic as
topological filters).
\end{theorem}

\begin{proof}
By Theorem~\ref{thm:sync}, decoder-phase-locking implies
$\mathcal{D}_i\circ\sigma=\mathcal{D}_j$ for some isomorphism $\sigma$ on
every input. The aperture implements the decoder; if decoders are isomorphic,
apertures have the same multipole order and field-configuration geometry,
hence are isomorphic as filters.
\end{proof}

\begin{corollary}[Aperture-mediated synchronisation]\label{cor:aper-mediated}
A necessary condition for phase-locked coordination is the existence of a
common aperture $\Aper^*$ to which every $\Aper_i$ maps via a
structure-preserving isomorphism.
\end{corollary}

\begin{theorem}[Memory-trace compatibility]\label{thm:mem-compat}
A globally phase-locked ensemble $E$ has agent memory traces $M_i$ that are
pairwise congruent: there exists $\eps_{\mathrm{mem}}>0$ such that
$|M_i-M_j|\le\eps_{\mathrm{mem}}$ for all $i,j$, where
$\eps_{\mathrm{mem}}$ depends only on historical decoherence excursions.
\end{theorem}

\begin{proof}
By Theorem~\ref{thm:sync}, phase-locked agents have identical
decoder--projection chains. Their $H^+$-fields are the same input stream
processed through identical decoders, hence
$H^+_i(t)=H^+_j(t)+\eta_{ij}(t)$ for noise $\eta_{ij}$ from
internal-state fluctuations. Integrating gives
$|M_i-M_j|\le\int|\eta_{ij}|\,\mathrm{d}t$, bounded by the decoherence
excursion magnitude.
\end{proof}

\begin{remark}[Shared memory as consequence, not prerequisite]\label{rem:mem-interp}
Theorem~\ref{thm:mem-compat} explains why synchronised teams remember the
same things: not because they share an external memory store, but because
phase-locked decoders process the contextual stream identically. Shared
memory of a joint task is a \emph{consequence} of synchronisation, not a
prerequisite for it. This is observable in synchronised musical ensembles,
surgical teams, military formations, and aligned multi-agent AI systems.
\end{remark}

\begin{figure*}[!htbp]
    \centering
    \includegraphics[width=\textwidth]{figures/synchronized-coordination/panel_5.png}
    \caption{\textbf{Aperture Sharing and Memory Compatibility.}
    (\textbf{A})~Aperture multipole patterns in the $(x,y)$ plane for monopole
    ($\ell=0$, isotropic), dipole ($\ell=1$, two lobes), and quadrupole
    ($\ell=2$, four lobes). Phase-lock requires shared multipole order
    (Theorem~\ref{thm:aper-share}). (\textbf{B})~Aperture-sharing matrix for
    a 5-agent phase-locked ensemble (all dipole apertures): uniformly green,
    confirming pairwise isomorphism. (\textbf{C})~Memory traces $M(t)$ for
    four phase-locked agents processing a common $H^+$-field with small
    per-agent noise ($\sigma=0.02$): nearly indistinguishable, confirming
    Theorem~\ref{thm:mem-compat}. (\textbf{D})~Three-dimensional synchronisation
    phase diagram over $(K,\sigma_\omega)\in[0.1,4]\times[0.1,2.0]$: the
    surface rises with $K$ (stronger coupling synchronises more) and falls with
    $\sigma_\omega$ (greater heterogeneity resists synchronisation); zero on the
    critical-coupling locus $K=\Kc^{\mathrm{ens}}=2\sigma_\omega/\pi$.}
    \label{fig:ch04-panel5-aperture}
\end{figure*}

%==========================================================================
\section{Numerical validation}
\label{sec:ch04-validation}
%==========================================================================

The theoretical results are validated across fifty experiments in ten
clusters, each comparing measured quantities to closed-form predictions.
All fifty experiments pass at machine precision; Table~\ref{tab:ch04-validation}
summarises the results.

\begin{table}[h!]\centering
\begin{tabular}{lcc}
\toprule
\textbf{Cluster (\# experiments)} & \textbf{Max rel.\ error} & \textbf{Verdict}\\
\midrule
C1 Single-agent state manifold (E1--E5)   & $0.00$            & PASS \\
C2 Partition potential and gradient flow (E6--E10) & $\le 10^{-15}$ & PASS \\
C3 Single-agent five-regime classification (E11--E15) & $0.00$ (boolean) & PASS \\
C4 Bridge to receiver agent (E16--E20)    & $\le 10^{-15}$    & PASS \\
C5 Trajectory-compatible coordination (E21--E25) & $0.00$ (boolean) & PASS \\
C6 Ensemble Kuramoto order parameter (E26--E30) & $\le 10^{-15}$ & PASS \\
C7 Five coordination regimes (E31--E35)   & $0.00$ (boolean)  & PASS \\
C8 Joint Lagrangian dynamics (E36--E40)   & $\le 10^{-14}$    & PASS \\
C9 Synchronisation as partition extinction (E41--E45) & $0.00$ (boolean) & PASS \\
C10 Aperture sharing and memory (E46--E50)& $\le 10^{-15}$    & PASS \\
\midrule
\textbf{Aggregate (50)} & $\mathbf{\le 10^{-14}}$ & \textbf{50/50 PASS} \\
\bottomrule
\end{tabular}
\caption{Validation summary across 50 experiments. All verdicts PASS at
machine precision.}
\label{tab:ch04-validation}
\end{table}

%==========================================================================
\section{Summary}
\label{sec:ch04-summary}
%==========================================================================

This chapter developed the Lagrangian agent calculus extending the thin
receiver--methodology--goal triple of Chapter~\ref{ch:purpose}. The
principal results are:

\begin{enumerate}[label=(\arabic*),leftmargin=*]
\item \emph{Three axioms and five-dimensional state manifold}
(Section~\ref{sec:ch04-axioms}--\ref{sec:ch04-agent}): bounded phase space,
no null state, and finite resolution underlie the compact manifold
$\Mfld=[0,1]\times[0,2\pi^2]\times[0,1]^3$, parametrised by Kuramoto order,
phase variance, and three S-entropy coordinates.
\item \emph{Partition potential and five operational regimes}
(Theorem~\ref{thm:five-regimes}): the Onsager--Machlup action generates
gradient flow on $\Mfld$; the partition potential's four terms produce a
Landau bifurcation at $K=\Kc$ that underlies all regime transitions. Five
regimes partition the $R$-axis at boundaries $\{0.3,0.5,0.8,0.95\}$, each a
second-order phase transition.
\item \emph{Forgetful projection and cell-truth bridge}
(Theorem~\ref{thm:proj-consist}, \ref{thm:ch04-cell-truth-rec}): the
Lagrangian agent contains the thin agent as a forgetful projection. Cell-truth,
Common-Cell Convergence, the catalytic composition law, motivation
heterogeneity, and the classical-impossibility reformulations of
Chapters~\ref{ch:subtask}--\ref{ch:purpose} are recovered as corollaries.
Trajectory-compatible coordination adds trajectory and memory compatibility
requirements to cell-coordination.
\item \emph{Five ensemble coordination regimes} (Theorem~\ref{thm:ens-five},
\ref{thm:joint-lag}): an ensemble inherits the five-regime classification
via the ensemble Kuramoto order parameter $\Ord$; transitions are second-order;
a single joint Lagrangian $\Lag_{\mathrm{ens}}=\sum_i\Lag^{(i)}+\Lag_{\mathrm{coup}}$
governs all regimes.
\item \emph{Synchronisation as partition extinction} (Theorem~\ref{thm:sync}):
the phase-locked regime is mathematically isomorphic to partition extinction
--- the same phenomenon underlying superconductivity and superfluidity. Global
phase-locking ($\Ord\ge 0.95$), zero inter-agent partition lag, and zero
coordination friction are equivalent conditions, and the transition is
discontinuous. The critical coupling is $\Kc^{\mathrm{ens}}=2\sigma_\omega/\pi$
and the transition is reversible (Corollary~\ref{cor:reverse}).
\item \emph{Aperture sharing and memory compatibility}
(Theorems~\ref{thm:aper-share}, \ref{thm:mem-compat}): phase-locked agents
share apertures and maintain congruent memory traces. Shared memory is a
\emph{consequence} of synchronisation, not a prerequisite.
\end{enumerate}

Chapters~\ref{ch:subtask}--\ref{ch:cohesion} together constitute the formal
Part~I of the monograph: a complete mathematical theory of bounded knowing,
bounded inquiry, bounded coordination, and cohesion dynamics. A single
bounded receiver cannot eliminate its floor; an ensemble drives its composite
floor to zero asymptotically; when agents phase-lock, inter-agent friction
vanishes and coordination becomes frictionless. Part~II descends from the
abstract calculus into its biological substrate: the thalamic gate as the
physical receiver, synaptic pharmacology as the modulator of floor noise, and
the dreaming brain as the extreme case --- full aperture closure, $\Sfloor\to 0$,
fabrication unconstrained.
